\label{chap.8}

\section{Summary of Research Work}

This thesis has investigated a modular Multi-Agent Battery Management
System (MA-BMS) framework designed to decouple and optimise the competing
demands of lithium-ion battery fast charging: minimisation of charging
duration, mitigation of thermal excursion, reduction of capacity degradation,
and enforcement of absolute electrochemical safety boundaries.
Unlike traditional monolithic or scalarised deep reinforcement learning
formulations, the proposed architecture separates charging speed, thermal
regulation, and health preservation into autonomous proposer agents, with
an analytical, non-learning safety supervisor acting as the deterministic
final authority.

The investigation was grounded in a comprehensive empirical audit of the
NASA Prognostics Center of Excellence (PCoE) Accelerated Life Testing (ALT)
dataset, comprising 26 two-cell series Samsung INR18650-25R packs cycled
across 4{,}918{,}193 raw telemetry rows. Key findings and contributions
of this thesis include:

\begin{enumerate}
  \item \textbf{Dataset Audit and Protocol Characterisation}:
  A rigorous block-level audit established that the NASA charging protocol
  is constant-current (CC) at \SI{3.0}{A} (1.2C) with a negligible
  constant-voltage (CV) plateau fraction ($\le 4\%$). The audit resolved
  ambiguities in the literature by separating three recommissioned packs
  ($\mathrm{baseline} < \SI{2.0}{Ah}$) from the 23 fresh packs, and
  identified that 10 of the 15 terminal failure events occurred while
  pack capacity was strictly above 0.80 SOH.

  \item \textbf{Calibrated Ageing-Aware Battery Simulator}:
  An equivalent circuit model (ECM) coupled with a lumped first-order
  thermal model was calibrated against empirical telemetry. A novel
  ageing-dependent series resistance law,
  $R_0(\mathrm{SOH}) = R_0^{\mathrm{new}}(1 + k_{R_0}(1-\mathrm{SOH}))$,
  was incorporated with stochastic coefficient sampling
  $k_{R_0} \sim \mathcal{U}(1.5, 4.5)$ based on cross-pack calibration
  ($r = 0.54$). Thermal response was validated against 5{,}479 physical
  rest cooling blocks ($\tau_{\mathrm{th}} = \SI{791}{s}$).

  \item \textbf{Four-State Adaptive State Estimation}:
  A four-state Extended Kalman Filter (EKF) tracking SOC, diffusion voltage
  $V_1$, capacity scaling $b$, and resistance scaling $\rho$ was developed
  with an ageing-aware initialisation: $\rho_0 = 1 + 3.0(1 - \mathrm{SOH}_0)$.
  The adaptive filter eliminated systematic voltage underprediction and
  reduced unsafe executed steps to zero across nominal and noisy episodes.

  \item \textbf{Rigorous Empirical Benchmarking}:
  Across nine unseen test episodes, a tuned single-agent policy reduced
  charging duration by 20.1\% relative to conventional CC-CV ($1{,}435\,\mathrm{s}$
  vs $1{,}796\,\mathrm{s}$) with zero safety violations. A 10{,}000-sample
  symmetric Dirichlet weight sensitivity analysis demonstrated that the
  single-agent controller achieved the lowest composite objective in
  84.78\% of random preference configurations.

  \item \textbf{Deterministic Safety Enforcement Under Faults}:
  Under seven injected sensor fault regimes (including stuck readings,
  noise dropouts, and calibration biases), the analytical supervisor
  completely eliminated hard voltage and temperature violations, reducing
  unsafe operation from 70--100\% (without supervisor) to 0.0\%.

  \item \textbf{Survival Analysis and LOPO RUL Prediction}:
  An interval-censored Kaplan--Meier analysis demonstrated that composite
  operational lifetime (Endpoint B: median \SI{630.5}{Ah}) is nearly
  19\% shorter than indicated by capacity fade alone (Endpoint A: median
  \SI{777.7}{Ah}). A Leave-One-Pack-Out (LOPO) Ridge regression model
  achieved an RUL MAE of \SI{113.0}{Ah} ($R^2 = 0.732$), outperforming
  population and linear baselines.
\end{enumerate}

\section{Resolution of Research Questions}

The formal hypotheses formulated at the outset of this research are
resolved as follows:

\subsection*{Research Question 1 (RQ1: Health-Cap Degradation Benefit)}
\begin{quote}
  \emph{Does an explicit health-cap proposer agent produce a statistically
  significant reduction in capacity degradation relative to an independent-deadline
  baseline?}
\end{quote}
\textbf{Status: Rejected / Not Supported.}
Empirical evaluation in the calibrated simulator revealed that the capacity
loss reduction achieved by the health-cap controller ($\Delta Q = 0.000143$)
compared to the independent baseline ($\Delta Q = 0.000151$) is negligible
and within the stochastic noise of the simulation. More importantly, when
the ageing resistance mechanism $R_0(\mathrm{SOH})$ was introduced, the
health-cap controller suffered severe deadline failure: its target success
rate dropped from 93.3\% to 65.6\% when initial SOH was sampled from
$\mathcal{U}(0.75, 0.80)$. Because internal resistance surges in aged cells,
the health cap triggers premature thermal throttling, massively prolonging
charge time without delivering meaningful electrochemical life extension.
Consequently, RQ1 is formally not supported under realistic aged cell dynamics.

\subsection*{Research Question 2 (RQ2: Multi-Agent Coordination Superiority)}
\begin{quote}
  \emph{Does dynamic arbitration between decoupled proposer agents outperform
  a tuned single-agent policy across diverse multi-objective weighting regimes?}
\end{quote}
\textbf{Status: Qualified / Open Horizon.}
In the static and Dirichlet weight benchmark, the Bayesian-tuned single-agent
policy dominated composite performance, securing 84.78\% best-weight share
due to its simultaneous minimisation of charging duration and cooling energy.
While all four controllers are Pareto-nondominated across the three-dimensional
objective surface (time, degradation, energy), dynamic arbitration between
competing agents requires a game-theoretic or priority-negotiating coordinator.
Without an active dynamic coordinator, the single-agent tuned policy remains
the most practical and high-performing controller.

\subsection*{Research Question 3 (RQ3: Deterministic Supervisor Robustness)}
\begin{quote}
  \emph{Can an analytical, non-learning safety supervisor enforce zero hard
  electrochemical and thermal violations under sensor corruption and fault regimes?}
\end{quote}
\textbf{Status: Strongly Supported.}
The deterministic safety supervisor achieved a 0\% hard violation rate across
all evaluated fault modes: voltage dropout (5\%), voltage positive bias
($+0.1\,\mathrm{V}$), temperature stuck, temperature dropout (5\%), and
moderate thermal disturbance. In contrast, unconstrained neural or heuristic
proposers suffered 70\% to 100\% hard violation rates, with up to 111 unsafe
seconds per episode. The supervisor successfully decouples performance
optimisation from safety preservation.

\section{Limitations of the Study}

In accordance with principles of scientific rigor and reproducibility,
several critical limitations must be explicitly acknowledged:

\begin{enumerate}
  \item \textbf{Sensor Bias Detection Horizon}:
  Under a constant negative temperature bias of $-5\,^\circ\mathrm{C}$
  (sensor permanently reads $5\,^\circ\mathrm{C}$ cooler than truth),
  the thermal plausibility monitor cannot detect the fault because the
  instantaneous thermal residual $|T_{\mathrm{meas}} - \hat{T}|$ remains
  bounded within the detection threshold. Consequently, a small residual
  of 4.2 unsafe executed steps per episode persisted under this specific fault.
  A single-channel residual monitor is mathematically incapable of distinguishing
  a constant sensor bias from unmodelled convective boundary variations.

  \item \textbf{Charge Telemetry Observability}:
  The NASA ALT dataset did not record charging current telemetry; charging
  current (\SI{3.0}{A}) is documented in the accompanying published papers
  and verified by integrating Ah between rests. Only the thermal cooling
  response during rest blocks is physically validated against raw charge data.
  Voltage dynamics and impedance during the charge phase are simulated from
  discharge-calibrated parameters rather than directly fitted to charge telemetry.

  \item \textbf{Absence of Early Failure Precursors}:
  Extensive statistical screening of heating amplitudes, thermal time
  constants, and voltage residuals revealed no detectable precursor signals
  prior to catastrophic electrical collapse. For the 10 packs that failed
  above 0.80 SOH, terminal collapse occurred abruptly, indicating that
  purely surface-level thermal and voltage monitoring cannot anticipate
  internal mechanical fracture or separator puncture.
\end{enumerate}

\section{Scope for Future Research}

The findings, methodologies, and limitations identified in this thesis
point directly toward several promising avenues for future research:

\begin{enumerate}
  \item \textbf{Multi-Sensor Analytical Redundancy and Voting}:
  To overcome the vulnerability of the single thermal monitor to constant
  sensor biases, future BMS designs should incorporate dual-thermistor
  configurations with a parity-space voting algorithm. Comparing differential
  spatial heat flux across cell terminals allows rapid isolation of sensor
  bias versus true electrochemical heating.

  \item \textbf{Dynamic Multi-Agent Coordination via Bargaining Games}:
  Rather than relying on static weight vectors, a higher-level coordinator
  can be formulated as a Nash bargaining game or constrained Markov decision
  process. Under this framework, the coordinator dynamically shifts bargaining
  power to the thermal agent when ambient temperatures surge, and to the
  charging speed agent when approaching critical delivery deadlines.

  \item \textbf{Hardware-in-the-Loop (HIL) Embedded Implementation}:
  The algorithms developed in this thesis are computationally lightweight:
  the 4-state adaptive EKF requires only $4 \times 4$ matrix operations,
  and the safety supervisor evaluates closed-form analytical equations.
  Deploying and benchmarking the compiled C/C++ firmware on an automotive-grade
  microcontroller (e.g., ARM Cortex-M7 or Infineon AURIX) connected to a
  real-time battery pack cycler is the logical next step toward commercialisation.

  \item \textbf{Physics-Informed Neural Network (PINN) State Observers}:
  While equivalent circuit models offer real-time execution, they cannot
  track lithium concentration gradients or solid-electrolyte interphase (SEI)
  film growth directly. Developing reduced-order PINNs that approximate the
  Doyle-Fuller-Newman electrochemical model in sub-millisecond execution times
  could enable direct observation of lithium plating overpotential.
\end{enumerate}

\section{Concluding Remarks}

As battery-powered systems assume an increasingly ubiquitous role in modern
infrastructure, the tension between fast charging performance, pack longevity,
and operational safety will only intensify. This thesis has demonstrated
that safety must not be left to empirical neural network training or soft
penalty tuning; instead, safety must be guaranteed deterministically by an
analytical supervisory boundary. By pairing an adaptive multi-state observer
with decoupled proposers and a hard safety projection, the Multi-Agent BMS
architecture achieves fast, energy-efficient charging while maintaining
bulletproof safety guarantees across the entire lifecycle of the battery.
